\documentclass[11pt]{article}
\usepackage[T1]{fontenc}
\usepackage{lmodern,amsmath,amssymb,amsthm,mathtools,microtype}
\usepackage[letterpaper,margin=.9in]{geometry}
\usepackage{booktabs,array}
\usepackage{hyperref}
\hypersetup{pdfauthor={Peyton Wang and Daniel Wang},pdftitle={No Fold Certificate of Darkness Below 21 Walls in Orthogonal Lattice Rooms},colorlinks=true,linkcolor=black,citecolor=black,urlcolor=blue}
\setlength{\parindent}{0pt}
\setlength{\parskip}{4pt plus 1pt minus 1pt}
\newtheorem{theorem}{Theorem}
\newtheorem{lemma}{Lemma}
\newtheorem{proposition}{Proposition}
\newtheorem{corollary}{Corollary}
\theoremstyle{definition}
\newtheorem{definition}{Definition}
\theoremstyle{remark}
\newtheorem{remark}{Remark}
\newcommand{\Q}{\mathbb Q}
\newcommand{\Z}{\mathbb Z}
\newcommand{\R}{\mathbb R}
\newcommand{\Par}{\mathrm{Par}}
\newcommand{\Rig}{\mathrm{Rig}}
\newcommand{\Cut}{\mathrm{Cut}}
\newcommand{\Rr}{R_{22}}
\newcommand{\src}[1]{\path{#1}}
\title{No Fold Certificate of Darkness Below 21 Walls\\in Orthogonal Lattice Rooms}
\author{Peyton Wang\\\small Francis Parker School \and Daniel Wang}
\date{Working draft v0.1 --- 23 November 2026\\\small Not submitted or peer reviewed}
\begin{document}
\maketitle
\begin{abstract}
A pair of interior points of a polygonal room is dark when no billiard path
that avoids the corners joins them. For orthogonal rooms with integer
vertices, the double of the room is a square-tiled cover of a pillowcase, and
three algebraic certificates (Par, Cut, and Rig) over a pillowcase base each
prove that a pair is dark at one target parity class. We prove, with computer
assistance, that no strictly simple orthogonal lattice room with at most
20 walls has a pair of rational interior points certified dark by these
certificates over any pillowcase base. The proof combines lemmas on
pillowcase bases with exhaustive checks of about $14.4$ billion rooms and at
least $59.4$ billion point pairs, which find no certified pair. A second proof
of the core case enumerates fold quotients and their covers instead of rooms.
The known 22-wall dark room serves as the positive control: its dark pair is
certified. We also record the fold-converse lemma, which describes exactly
when two points with equal images in a factor are joined, and a later result
that an R22-type parity fold needs at least nine reflex corners. The theorem
shows only that this certificate mechanism produces no dark pair below
21 walls. It does not show that every pair is lit.
\end{abstract}

\section{Statement and conventions}\label{sec:statement}

\paragraph{Rooms.}
A \emph{room} $P$ is a finite union of closed unit lattice squares (cells)
whose interior is connected, with no hole and no \emph{pinch}: a lattice point
with exactly two diagonally opposite cells. Its boundary is then one simple
closed curve, and we call $P$ \emph{strictly simple}. The wall count $W$ is the
number of maximal straight walls; it equals the number of corners. With $r$
reflex corners, $W=2r+4$, so $W\le 20$ means $r\le 8$. The \emph{area} $n$ is
the number of cells. The room is \emph{normalized} when its horizontal edge
lengths have gcd $1$ and its vertical edge lengths have gcd $1$. A lattice
point is a convex corner, a smooth wall point, a reflex corner, or an interior
point when $1$, $2$, $3$, or $4$ cells meet there.

\paragraph{Light.}
A \emph{lighting path} for distinct points $S,T$ is a billiard path from $S$ to
$T$ that reflects specularly off walls and never meets a corner. Corner
avoidance applies to every corner, convex corners included. The pair is
\emph{lit} when such a path exists and \emph{dark} otherwise.

\paragraph{Main result.}
Section~\ref{sec:certificates} defines pillowcase bases $P_\Lambda$, the
three certificates, and what it means for a pair to be \emph{certified dark}.

\begin{theorem}[Fold side at $W\le 20$; computer-assisted]\label{thm:main}
Let $P$ be a strictly simple orthogonal lattice room with $W\le 20$. Then no
pair of distinct rational interior points of $P$ is certified dark. That is,
for no admissible lattice $\Lambda$ do all four lift pairs of the pair have,
at every target, a $\Par$, $\Cut$, or $\Rig$ certificate over the pillowcase
base $P_\Lambda$.
\end{theorem}

The 22-wall room $\Rr$ of \cite{Q22} has $r=9$ and one dark pair, and that
pair is certified dark by $\Par$ at all three targets
(Section~\ref{sec:controls}). So the wall bound in Theorem~\ref{thm:main}
can't be raised to $22$.

\begin{remark}[What the theorem does not say]\label{rem:scope}
Theorem~\ref{thm:main} is a statement about certificates. It shows that this
certificate mechanism produces no dark pair below 21 walls. It does not show
that every pair is lit. A pair that carries no certificate may still be dark.
The statement that every dark pair is certified is the open conjecture $U^*$
of Section~\ref{sec:converse}; its per-class form is the realization lemma
there. Theorem~\ref{thm:main} also says nothing about irrational pairs
directly, because a pillowcase base exists only for rational pairs
(Section~\ref{sec:bases}).
\end{remark}

\paragraph{Sources and status.}
Every lemma below is taken from a numbered research memo or referee report in
the public repository. Section~\ref{sec:repro} gives the path convention.
Each theory input has had one independent AI referee and no human referee.
Where a referee repaired a step, we state the repaired version; those repairs
have not been reviewed again. The finite parts are computations. We mark each
result as \emph{proved}, \emph{computed}, or \emph{computer-assisted} (a proof
whose finite part is a computation).

\paragraph{Organization.}
Section~\ref{sec:double} sets up the double as a cover of the pillowcase.
Section~\ref{sec:certificates} defines certificates and pillowcase bases.
Section~\ref{sec:proof} proves Theorem~\ref{thm:main} from its lemmas and the
computations of Section~\ref{sec:computation}. Section~\ref{sec:second} gives
the second proof by fold enumeration. Section~\ref{sec:converse} states the
fold-converse lemma. Section~\ref{sec:r9} records the later result that an
R22-type parity fold needs $r\ge 9$. Section~\ref{sec:limits} lists the
limitations, and Section~\ref{sec:repro} explains how to reproduce the
computations.

\section{The double as a cover of the pillowcase}\label{sec:double}

\paragraph{The pillowcase and the double.}
Let $B$ be the unit pillowcase: two unit squares glued along their boundary.
Its four corners are the classes of $\Z^2/2\Z^2$. The double $DP$ is two
copies of $P$ glued along the boundary. It is an $n$-sheeted square-tiled
cover $DP\to B$, branched over the four corners, whose sheets are the cells of
$P$. The seam reflection $\iota$ of $DP$ exchanges the two copies.

\paragraph{Monodromy.}
Each corner $x$ of $B$ has a monodromy permutation $g_x$ of the cells. It
rotates by $\pi$ about the class-$x$ lattice points, and its cycles are the
stars of those points:
a $1$-cycle (\emph{pole}) is a convex corner;
a $2$-cycle (\emph{regular point}) is a smooth wall point or one half of an
interior point;
a $3$-cycle (\emph{zero}) is a reflex corner.
An interior point gives two $2$-cycles, so it has two \emph{lifts} in $DP$, and
a pair $(S,T)$ has four \emph{lift pairs}. The surface $DP$ has genus $0$. Room
corners are exactly the non-regular points of $DP$, so corner avoidance means
that a light path passes only through regular points.

\paragraph{Parity classes and targets.}
Translate so that the class of $S$ is $a=(0,0)$, and call the other three
classes $b$, $c$, $d$. A lighting path from $S$ to $T$ unfolds to a segment of
primitive slope $(p,q)$. Its image in $B$ runs back and forth along one saddle
connection from $a$ to the corner $x$ given by $(p,q)\bmod 2$. This $x$ is the
\emph{target} of the path. For a same-class pair the targets are the three
other classes; for a different-class pair the only target is the class of $T$.

\paragraph{Factors.}
A \emph{factor} is a quotient $\phi\colon Q\to N=Q/\Pi$ by a block system $\Pi$
of the monodromy group $G$. Here $Q$ is $DP$ or an intermediate cover. $N$ is
again a genus-$0$ pillowcase cover. Write $\delta=\deg(Q\to N)$,
$d=\deg(N\to B)$, so $n=\delta d$, and let $z$ and $p$ count the zeros and
poles of $N$. Gauss--Bonnet gives $p=z+4$.

\begin{lemma}[Fact L; d10, reviewed]\label{lem:factL}
Let $q\in Q$ lie over $u=\phi(q)$ with local degree $e$. Then the cycle length
of $q$ is $e$ times that of $u$. So over a regular point of $N$ every point of
$Q$ is regular with $e=1$; over a zero every point is a zero with $e=1$; over a
pole a point is a pole, a regular point ($e=2$), or a zero ($e=3$). In
particular $\phi$ is branched only over poles of $N$.
\end{lemma}

A point of $N$ is \emph{live} when a regular point of $Q$ lies over it.
Regular points of $N$ are live, zeros never are, and a pole is live when it has
a preimage with $e=2$. For a simple arc $\alpha$ from $a$ to a corner $x$, the
graph $\Gamma_\alpha(Q)$ is the bipartite graph of lifts of $\alpha$. The
\emph{live component} $K_N(\alpha,u)$ is the component of $u$ in the subgraph of
$\Gamma_\alpha(N)$ induced on live vertices.

\begin{lemma}[Projection; d10]\label{lem:proj}
$\phi$ maps the component of $\sigma$ in the regular subgraph of
$\Gamma_\alpha(Q)$ into $K_N(\alpha,\phi(\sigma))$.
\end{lemma}
\begin{proof}
$\phi$ maps lifts of $\alpha$ to lifts of $\alpha$, so it is a graph map.
Regular vertices go to live vertices, and the image of a connected set is
connected.
\end{proof}

\section{Certificates}\label{sec:certificates}

\subsection{Pillowcase bases}\label{sec:bases}
For a lattice $\Lambda\subset\R^2$ put
$P_\Lambda=\R^2/\langle 2\Lambda,-1\rangle$. Its corners are the images of
$\Lambda$, in the four classes $\Lambda/2\Lambda$. Write
$m=[\Lambda:\Z^2]$ for the \emph{index} and
$\Lambda_r=(\Lambda\cap\Q e_1)+(\Lambda\cap\Q e_2)$ for the largest
rectangular sublattice. Let $\rho(x,y)=(x,-y)$.

\begin{proposition}[Classification of bases; d47 Prop.~1.1, checked in d48]\label{prop:bases}
Let $R$ be a normalized strictly simple orthogonal lattice polygon with a
vertex at the origin. The pillowcases that $DR$ covers by a half-translation
map are, up to isometry, exactly the $P_\Lambda$ with $\Lambda\supseteq\Z^2$.
For a normalized room, the only pillowcase factor of $DR\to B$ is $B$ itself
(d10 R1, corrected in the d10--d12 review).
\end{proposition}

\begin{proof}[Proof summary]
Normalize the front chart to the identity; the back chart is $\rho$ up to a
translation. Gluing along a horizontal or vertical edge forces $(0,y_0)$ and
$(x_0,0)$ into $\Lambda$ for every vertex coordinate. Normalization makes these
generate $\Z^2$. Conversely the identity and $\rho$ glue correctly for every
$\Lambda\supseteq\Z^2$.
\end{proof}

Every lattice room is $P=\operatorname{diag}(g,g')R$ plus a translation for a
unique normalized $R$, where $g,g'$ are the axis gcds. $DP$ and $DR$ are the
same flat surface up to this affine map, and straight segments and cone points
correspond (d47 \S1). So we may work with $R$.

\begin{definition}[Admissible base; d47 Def.~1.2 with the d48 correction]\label{def:admissible}
Let $S\ne T$ be interior points of $R$. A lattice $\Lambda$ is
\emph{admissible} for $(S,T)$ when all four lifts lie over corners of
$P_\Lambda$, that is, $S,T,\rho S,\rho T\in\Lambda$. An admissible lattice
exists if and only if $S$ and $T$ are rational.
\end{definition}

The d48 referee found that d47's original condition
$\Lambda\supseteq\Z^2+\Z S+\Z T$ places only the front lifts over corners;
the condition above is the corrected one.

\subsection{The three certificates}
Fix a pillowcase base $P'$ (for example $P_\Lambda$) such that the lifts
$\sigma$ of $S$ and $\tau$ of $T$ lie over corners $c_s$ and $c_t$ of $P'$.
The targets are the three corners other than $c_s$ when $c_s=c_t$, and $c_t$
alone otherwise. A factor $\phi\colon Q\to N\to P'$ is a \emph{certificate at
the target $x$} when one of the following holds (d10 \S2.1).
\begin{itemize}
\item[$\Par(x)$:] $c_s=c_t$, $\phi(\sigma)\ne\phi(\tau)$, and $N$ has no
regular point over $x$.
\item[$\Rig(x)$:] $N\to P'$ has identity monodromy at a corner $y\notin\{c_s,x\}$,
and $\phi(\tau)\notin K_N(\alpha,\phi(\sigma))$ for one arc $\alpha$ from $c_s$
to $x$.
\item[$\Cut(x)$:] $c_s=c_t$ and $\phi(\sigma)=\phi(\tau)$; every regular point
of $Q$ over $c_s$ and over $x$ maps to a regular point of $N$; and the
monodromy of $N$ at a corner outside $\{c_s,x\}$ is a single cycle.
\end{itemize}

\begin{definition}[Certified dark]\label{def:certified}
A target is \emph{certified} for a lift pair when some factor is a certificate
there; the factor may change with the target. A rational pair $(S,T)$ is
\emph{certified dark} when, for some admissible $\Lambda$, every target of every
one of the four lift pairs is certified over $P_\Lambda$.
\end{definition}

\begin{theorem}[Soundness, Theorem U; d10, reviewed]\label{thm:U}
If some pillowcase base has a certificate at every target, the pair is dark.
\end{theorem}

\begin{proof}[Proof (d10 \S2.2)]
A light path projects to $P'$ as a straight path from $c_s$ that runs back and
forth along one saddle connection to a target $x$. It suffices to show that
$\tau$ is outside the regular component of $\sigma$ in $\Gamma_\alpha(Q)$ for
every arc $\alpha$ from $c_s$ to $x$.
\emph{Par.} A live vertex of $\Gamma_\alpha(N)$ over $x$ is a pole, hence a
leaf. So $\phi(\sigma)$ is the only vertex of its live component over $c_s$,
and Lemma~\ref{lem:proj} finishes.
\emph{Rig.} $N\to P'$ is unbranched over $y$, so $N$ covers the sphere minus
the other three corners. All simple arcs from $c_s$ to $x$ are isotopic there
and the isotopy lifts, so one arc decides all arcs.
\emph{Cut.} Let $C$ be the component of $\phi(\sigma)$ in the regular subgraph
of $\Gamma_\alpha(N)$. Every vertex of the light component maps to a regular
point, so $\phi$ is injective on it when $C$ is a path. When $C$ is a cycle,
the arc $\beta$ between the other two corners has connected preimage disjoint
from $C$, and one disc bounded by $C$ contains no pole. By
Lemma~\ref{lem:factL}, $\phi$ is unbranched over that disc, the lift of $C$ is
closed, and $\phi$ is again injective. In both cases $\tau\ne\sigma$ cannot share
the image of $\sigma$ in the component.
\end{proof}

The proof never uses $P'=B$, so it applies over every admissible base
(d47 \S1, remark after Def.~1.2). A $\Cut$ certificate has
$\deg(N\to P')\le 3$ (d10; reworded in d50).

\subsection{Canonical factors}
For a same-class pair over $B$, let $B_x$ be the finest block system that
merges every $2$-cycle of $g_x$, and $N_x=DP/B_x$. Let $R_y$ be the finest
block system that merges every cycle of $g_y$.

\begin{proposition}[R3; d10, reviewed]\label{prop:R3}
(1) $N_x$ has no regular point over $x$, and $\Par(x)$ holds for some factor
exactly when $B_x$ separates the lift pair, that is, when the two lifts have
different images in $N_x$.
(2) $\Rig(x)$ with identity at $y$ holds for some factor exactly when it holds
over $DP/R_y$.
(3) $B_x$ and $R_y$ are invariant under $\iota$, so $N_x$ is the double of a
fold target $K_x=N_x/\iota$ of $P$.
\end{proposition}

\begin{remark}[Class types]
$B_x$ is one block, so $N_x=B$, or $B_x$ is discrete, so $N_x=DP$, or
$1<\deg N_x<n$. The second case occurs exactly when class $x$ has no interior
or smooth wall point (d68 Lemma B, as recorded in the results record).
\end{remark}

\section{Proof of the main theorem}\label{sec:proof}

The proof has five parts, following the assembly in the results record
(\src{d70-writeup/RESULTS.md}, \S3.3) and the coordinator's record of
17 November 2026 (\src{coordinator-log-archive-1.md}).

\begin{center}
\small
\begin{tabular}{@{}lp{0.5\textwidth}ll@{}}
\toprule
Part & Content & Status & Source \\
\midrule
F1 & Reduction of a certified same-class pair to three $\Par$ separations & proved (reviewed) & d39, d43 \\
F2 & Filter $F^*$ on area and reflex classes & proved (reviewed) & d44, d45 \\
F3 & Room enumeration at $W=20$, areas allowed by $F^*$ & computed & d46 \\
F4 & Pillowcase bases: odd index, even index, different-class marks & proved (reviewed once each) & d47--d50 \\
F5 & The two residual finite checks (Sets A and E) & computed & d56 \\
\bottomrule
\end{tabular}
\end{center}

\subsection{Lemmas on pillowcase bases (d47 to d50)}

\begin{lemma}[Corner classes; d47 Lemmas 1.3 and 1.4]\label{lem:classes}
Let $g,g'$ be the axis gcds of a lattice room. Its corners occupy at least
$3$ parity classes if and only if $g$ and $g'$ are both odd. Over a base
$P_\Lambda$ with $\Lambda\supseteq\Z^2$, the vertices of $R$ occupy at least
$3$ classes of $\Lambda/2\Lambda$ if and only if $m$ is odd; then
$\Z^2/2\Z^2\to\Lambda/2\Lambda$ is a bijection.
\end{lemma}

\begin{proposition}[Rectangular reduction; d47 Prop.~2.1]\label{prop:rect}
Let $\Lambda=\frac1a\Z\times\frac1b\Z$ be admissible and
$A=\operatorname{diag}(a,b)$. Then $R_{a,b}=AR$ is a strictly simple orthogonal
lattice room with the same $W$ and $r$ and area $ab\,n_R$, the points $AS,AT$
are lattice points, and a certificate system over $P_\Lambda$ is the same thing
as an ordinary certificate system over the unit pillowcase for
$(R_{a,b},AS,AT)$.
\end{proposition}

\begin{lemma}[Transfer; d47 Lemma 3.2 with the d48 fixes]\label{lem:transfer}
Let $\Lambda'\subseteq\Lambda$ be lattices with $[\Lambda:\Lambda']$ odd, such
that all four lifts lie in $\Lambda'$. Then every full certificate system over
$P_\Lambda$ yields one over $P_{\Lambda'}$.
\end{lemma}

The d47 proof is a sketch in five steps (corners, factor, $\Par$, $\Cut$,
$\Rig$). The d48 referee re-derived each step: the factor is the meet of the
two block systems; the $\Rig$ arc is a saddle connection from a point of class
$c_s$ along a vector primitive in both lattices; and the cyclic hypothesis is
removed by composing prime-index steps. The d48 referee also proved that, for
odd $[\Lambda:\Lambda_r]$, $\Lambda\cap\rho\Lambda=\Lambda_r$. So for odd index
an admissible pair has $S,T\in\Lambda_r$, and Lemma~\ref{lem:transfer} with
$\Lambda'=\Lambda_r$ and Proposition~\ref{prop:rect} give:

\begin{corollary}[Odd index; d47 Cor.~3.3, d48]\label{cor:odd}
If $m$ is odd, a certificate system over $P_\Lambda$ reduces to an ordinary
certificate system over the unit pillowcase of $R_{a,b}$, where
$\Lambda_r=\frac1a\Z\times\frac1b\Z$ with $a,b$ odd. The room $R_{a,b}$ has
corners in at least $3$ classes, and same-class marks stay same-class
(d47 \S2).
\end{corollary}

For even index, write $\Lambda_{(2)}$ for the preimage of the $2$-primary part
of $\Lambda/\Z^2$, $M_0=\Z^2+\Z S+\Z T$, and $V$ for the image of $\Z^2$ in
$\Lambda/2\Lambda$. Room vertices lie over $V$, and $|V|\le 2$ for even $m$.

\begin{lemma}[Odd-part reduction; d48, d49 \S2, d50 fix]\label{lem:oddpart}
A certificate over $P_\Lambda$ transfers to one over $P_{\Lambda'}$ with
$\Lambda'=\Lambda_{(2)}+M_0+\rho M_0$. For lattice marks
$[\Lambda':\Z^2]$ is a power of $2$. For non-dyadic marks it is $2^k$ times an
odd number at least $3$, so $[\Lambda':\Z^2]\ge 6$.
\end{lemma}

The term $\rho M_0$ is the d50 referee's fix; without it $\rho S$ can fall
outside $\Lambda'$.

\begin{lemma}[Lemma Z; d49, reviewed by d50]\label{lem:Z}
Over any base $P_\Lambda$, let a certificate have degree $n=\delta d$ with $z$
zeros. If $z\ge 1$, then $n\le 8r$ for $\Par$, $n\le 5r$ for $\Rig$, and
$n\le 3r$ for $\Cut$. A zero-free $\Cut$ has $d=2$, forces $n\le 4r+4$, and its
target class is vertex-only.
\end{lemma}

\begin{proof}[Proof summary]
Put $w_j=p_j+3z_j$ for each base class $j$. Then
$\sum_j w_j=p+3z=4z+4$. For $\Par$ at $x$, $d=w_x\le 4z+4\le 8z$, and each
zero of $N$ has $\delta$ zero preimages, so $\delta z\le r$ and $n\le 8r$. For
$\Rig$ with identity at $y$, $d=p_y\le z+4\le 5z$. For $\Cut$, a single cycle
has length at most $3$ (Lemma~\ref{lem:factL}), so $d\le 3$; $d=3$ puts a zero
at the cut class, and $d=2$ is zero-free with no regular point of $Q$ over $x$,
so every point over $x$ is a vertex. The d50 referee corrected d49's wording
``$\Cut$ forces $d=3$'', which holds only with a zero.
\end{proof}

\begin{lemma}[Lemma Z$'$; d49, reviewed by d50]\label{lem:Zprime}
At a target outside $V$, any certificate with a zero has $n\le 4r$; the same
bound holds for any $\Rig$ certificate with a zero.
\end{lemma}

\begin{lemma}[Lemma R0; d49, reviewed by d50]\label{lem:R0}
A zero-free $\Rig$ needs $d=4$, $\Lambda\supseteq\frac12\Z^2$, and every vertex
regular. So no zero-free $\Rig$ exists for lattice marks or at $m=2$.
\end{lemma}

\begin{lemma}[Witness lemma; d49, reviewed by d50]\label{lem:witness}
Zero-free $\Par$ and zero-free $\Rig$ certificates depend only on $\Lambda$,
the factor lattice, the lift pair, and the classes of $S,T$ modulo $2\Z^2$.
For any such data, the L-hexagon $[-2K,2K+1]^2\setminus[2K,2K+1]^2$, translated
to put a vertex at the origin, lights the lift pair. So no certificate that
depends only on lattice data separates all targets of a lift pair.
\end{lemma}

The per-target version of Lemma~\ref{lem:witness} is false: for
$\Lambda=\frac12\Z\times\Z$ and lattice marks with odd $T_x-S_x$, the factor
$B$ separates at the target in $V$ in every room (d49 \S4). The d50 referee
fixed the wording ``classes mod $2\Z^2$'' and noted that zero-free $\Cut$ is
room data, not lattice data.

\begin{proposition}[Even-index bounds; d49 \S5 with the d50 repairs]\label{prop:even}
Let $m'=[\Lambda':\Z^2]$ be even as in Lemma~\ref{lem:oddpart}, and suppose all
four lift pairs are certified over $P_{\Lambda'}$. Then $n_R\le 2r\le 16$.
\end{proposition}

\begin{proof}[Proof summary]
At $m'=2$, a zero-free factor is $P_{\Lambda''}$ with $d\mid m'$. For $d=1$ no
certificate exists; for $d=2$, $N=B_R$ has one regular point over each non-$V$
corner and no vertex over it, and $\Rig$ is excluded by Lemma~\ref{lem:R0}.
Every lift pair with $c_s=c_t$ has a non-$V$ target, since $|V|\le 2$. So a
certificate there has a zero, and Lemma~\ref{lem:Zprime} gives $2n_R\le 4r$.
For $c_s\ne c_t$ only $\Rig$ applies, and Lemma~\ref{lem:R0} forces a zero. At
$m'\ge 4$, Lemma~\ref{lem:witness} forces some target to carry a certificate
with a zero or a zero-free $\Cut$, so $n\le 8r$ by Lemma~\ref{lem:Z} and
$n_R\le 8r/m'\le 2r$. This $m'=2$ argument is the d50 referee's repair; d49's
own citation did not cover $\Par$.
\end{proof}

\subsection{Different-class marks}

\begin{proposition}[d47 Prop.~4.1, corrected in d48]\label{prop:diffclass}
Let $R'$ have corners in at least $3$ classes and $r\le 8$, and let $S,T$ be
interior lattice points in different classes. A certificate for a lift pair
over the unit pillowcase exists only if $n\le 16$, or $n=20$ with $r=8$ and
$(\delta,d,z)=(4,5,2)$.
\end{proposition}

\begin{proof}[Proof summary]
With $c_s\ne c_t$ only $\Rig$ can certify, over $N=Q/R_y$ by
Proposition~\ref{prop:R3}. For $z=0$, $N$ is a pillowcase factor, so $N=B$ and
the single arc lift joins the two live corners. For $z=1$, $d\le 5$; at $d=3$
the arc graph is the path pole--regular--regular--pole with every vertex live,
and $d=4,5$ fail by cycle-type arguments that do not involve the marks. For
$z\ge 2$, $\delta z\le r$, $d\le z+4$, $\delta$ is even by the off-seam-pole
step of d39, and $d=6$ is excluded by the seam argument of d43. This leaves
$(\delta,z)=(2,\le 4)$ with $n\le 16$, and $(4,2)$ with $r=8$, $n=4d$. The d48
referee corrected d47's ``$d\in\{3,4,5\}$'' to $d=5$.
\end{proof}

\subsection{Same-class marks over the unit pillowcase (F1, F2)}

\begin{proposition}[F1; d39 Claims 2 to 4, reviewed in d43]\label{prop:F1}
Let $P$ be a strictly simple orthogonal lattice room with $r\le 8$ whose
corners occupy at least $3$ parity classes, and let $S\ne T$ be interior
lattice points of one class. Suppose every lift pair has, at each of $b,c,d$, a
$\Par$, $\Cut$, or $\Rig$ certificate over the unit pillowcase. Then $B_b$,
$B_c$, $B_d$ each separate all four lift pairs, and $5\le r\le 8$ and
$21\le n\le 5r\le 40$.
\end{proposition}

\begin{proof}[Proof summary]
\emph{Cut.} If class $x$ has no transparent point, $B_x$ is discrete and
$\Par$ already holds. Otherwise $\Cut$ forces $d\le 3$; $d=1$ has no regular
$a$ point, $d=2$ gives genus $1$ by Riemann--Hurwitz, and at $d=3$ the
patched R4 argument of the d10--d12 review forces the 12-wall cross, whose
pair is lit. So genuine $\Cut$ never appears in a full triple.
\emph{Essential Rig.} Over $N=DP/R_y$, zeros of $N$ lie on the seam, off-seam
poles come in at most one conjugate pair, and a $\Rig$ target has exactly one
interior $\pi$-cone, so $\delta$ is even. A target with $z=1$ gives no
essential $\Rig$. With $\delta z\le r$, $d\le z+4$, $z\ge 2$, and the exclusion
of $d=6$ (the regular $S_3$ cover, whose order-$3$ deck group puts at least
four poles off the seam), one gets $n\le 20$. The d43 referee filled two gaps:
the domino and square cases of the one-star step, and the proof that the
degree-$6$ cover is regular.
\emph{Census.} The d9 census finds no dark same-class interior pair at
$W\le 21$ and area at most $20$. By Theorem~\ref{thm:U}, a certified pair is
dark, so $n\ge 21$. Then every target is certified by $\Par$, which holds
exactly when $B_x$ separates (Proposition~\ref{prop:R3}). The d36.4 bound
$n\le 5r$ for a separating class quotient gives $r\ge 5$.
\end{proof}

\begin{proposition}[Filter $F^*$; d44, reviewed by d45]\label{prop:Fstar}
In the situation of Proposition~\ref{prop:F1}, let $\rho_y$ count reflex
corners in parity class $y$. Some factorization $n=\delta d$ with
$\delta\ge 3$ satisfies
\[
\sum_y\left\lfloor\rho_y/\delta\right\rfloor\ge z_{\min}(d),\qquad
z_{\min}(d)\,\delta\le r,
\]
where $z_{\min}(d)=\lceil (d-1)/4\rceil$ for odd $d$ and $\lceil d/4\rceil$ for
even $d$.
\end{proposition}

\begin{proof}[Proof summary]
If all three class quotients were Galois, d36.3 would make $B_b,B_c,B_d$ a
chain, so their join separates; the parity-factor theorem of d2 (claim 4)
excludes a common separating parity factor at $r\le 8$. So some target class
$x$ has a non-Galois quotient, hence $\delta\ge 3$. Each zero of $N_x$ has
$\delta$ reflex preimages in its own parity class, so $z\le\sum_y\lfloor
\rho_y/\delta\rfloor$, and d36.4 gives $z\ge z_{\min}(d)$. The d45 referee
checked that d2 applies to the d36.3 join and recomputed the area lists.
\end{proof}

At $r\le 8$, $F^*$ forces $z\le 2$ and $d\le 9$, and the allowed areas are:

\begin{center}
\small
\begin{tabular}{@{}cl@{}}
\toprule
$r$ & Areas allowed by $F^*$ in $21\le n\le 5r$ \\
\midrule
5 & 25 \\
6 & 21, 24, 25, 27, 30 \\
7 & 21, 24, 25, 27, 28, 30, 35 \\
8 & 21, 24, 25, 27, 28, 30, 32, 35, 36, 40 \\
\bottomrule
\end{tabular}
\end{center}

\subsection{Assembly}

\begin{proof}[Proof of Theorem~\ref{thm:main}]
Let $(S,T)$ be rational and suppose it is certified dark over an admissible
$\Lambda$. Pass to the normalized room $R$; then $\Lambda\supseteq\Z^2$ by
Proposition~\ref{prop:bases}, and $W$ and $r\le 8$ are unchanged.

\emph{Odd index.} By Corollary~\ref{cor:odd} the system is an ordinary
certificate system for lattice marks of $R_{a,b}$, which has the same $W$ and
corners in at least $3$ classes. If the marks are in one class,
Proposition~\ref{prop:F1} puts $R_{a,b}$ in the range $5\le r\le 8$,
$21\le n\le 5r$, with all three $B_x$ separating all four lift pairs.
Proposition~\ref{prop:Fstar} restricts $n$ to the listed areas. The d46 runs
(Section~\ref{sec:computation}) find no such room and pair. The d46 enumeration
covers all widths and heights, so non-normalized rooms with odd gcds are
included (d47 \S4a). If the marks are in different classes,
Proposition~\ref{prop:diffclass} gives $n\le 16$ or the residue D20 ($W=20$,
area $20$). For $n\le 16$, the d9 Part~A census lights every different-class
interior lattice pair at area at most $16$; for D20, the d56 Set~A run lights
every pair. An odd rescaling with $ab>1$ cannot reach D20, because it would need
a $W=20$ room of area $4$ (d47 \S4b). A lit pair has a lit lift pair, which
carries no full certificate by Theorem~\ref{thm:U}. This contradicts the
assumption.

\emph{Even index.} By Lemma~\ref{lem:oddpart} the system transfers to
$\Lambda'$, and Proposition~\ref{prop:even} gives $n_R\le 2r\le 16$. For
lattice marks, the d9 census (same class, area at most $20$) and Part~A
(different class, area at most $16$) light the pair, a contradiction (d49
\S6, d50 \S5). For non-lattice marks, the index of $\Lambda'$ is $2$ or $4$, or
$6$ for non-dyadic marks, which need $n_R\le\lfloor 8r/6\rfloor\le 10$ and
leave one $W=16$, area-$8$ room. Rooms in Part~A's one-cylinder classes (i) and
(ii) and rooms with $W\le 8$ have all pairs lit by d9 Part~A Task~2 and the
eight-wall theorem (d15). The remaining normalized class (iii) rooms with
$n_R\le 2r$ number $945{,}884$ (d50 \S7). The d56 Set~E run lights every
admissible non-lattice pair in them, a contradiction.
\end{proof}

\begin{remark}[Other bases]
The d48 referee also showed, as a sketch, that $\Par$ and $\Rig$ certificates
over a $(2,4,4)$ orbifold base reduce to pillowcase certificates, and that no
$\Cut$ is defined there. So Theorem~\ref{thm:main} needs no extra case for
those bases. The d51 memo, which is not reviewed, finds no comparable bound for
rooms with $45$-degree walls.
\end{remark}

\begin{remark}[Bases chosen per lift pair]
Definition~\ref{def:certified} uses one base for all four lift pairs. The d48
and d50 referees state that the conclusion also holds when each lift pair
chooses its own base, for lattice marks, except D20 when every lift pair uses
$\Z^2$; D20 is closed by Set~A. The d50 referee closed the remaining odd-index
per-lift-pair residue for marks outside $\Lambda_r$ by the d9 Part~A grids.
We state Theorem~\ref{thm:main} for one base.
\end{remark}

\section{The computation}\label{sec:computation}

\subsection{The room check (d46)}
The d46 program \src{d46-cert-check/cert_check.c} enumerates compressed
shapes (no two adjacent rows or columns equal) with positive column widths and
row heights. It applies $F^*$ per parity pattern and per area before any room
is built, then materializes each surviving room and tests whether some
same-class pair of interior lattice points has all four lift pairs separated
by each of $B_b$, $B_c$, $B_d$. Such a room and pair is a \emph{positive}.
Rooms with at most $2$ corner classes are skipped; by Lemma~\ref{lem:classes}
and Proposition~\ref{prop:bases} they are handled through their normal forms.
Each free room is generated once, up to the $8$ lattice symmetries. The design
and its sizing are in \src{d44-cert-check-sizing/README.md}.

\begin{table}[h]
\centering
\small
\begin{tabular}{@{}rrrrr@{}}
\toprule
Area & Rooms checked (free) & Fixed rooms & Pairs & Positives \\
\midrule
21 & 105,347,581 & 842,729,700 & 52,924,264 & 0 \\
24 & 961,558,778 & 7,692,288,677 & 1,013,958,816 & 0 \\
25 & 219,938,227 & 1,759,502,724 & 362,992,628 & 0 \\
27 & 6,106,330,113 & 48,850,258,466 & 11,508,487,504 & 0 \\
28 & 99,057,215 & 792,340,041 & 272,682,085 & 0 \\
30 & 289,576,463 & 2,316,525,088 & 1,231,691,360 & 0 \\
32 & 787,331,513 & 6,298,310,531 & 3,669,618,239 & 0 \\
35 & 123,433,861 & 987,469,564 & 1,089,342,824 & 0 \\
36 & 4,578,115,956 & 36,624,090,676 & 33,277,901,167 & 0 \\
40 & 15,380,590 & 123,017,620 & 246,266,377 & 0 \\
\midrule
Total & 13,286,070,297 & & 52,725,865,264 & 0 \\
\bottomrule
\end{tabular}
\caption{The d46 full run at $W=20$ ($r=8$), by area. There are
$1{,}058{,}080{,}070$ canonical shapes; $500$ jobs; $87{,}829.7$ CPU seconds.
Source: \protect\src{d46-cert-check/runs/full_W20/aggregate.md}.}
\label{tab:d46}
\end{table}

\subsection{The residual checks (d56)}
The d56 checker \src{d56-final-check/src/check.c} handles the two residues
left by Section~\ref{sec:proof}. Set~A is D20: all strictly simple rooms with
$W=20$ and area $20$, with different-class interior lattice pairs over $\Z^2$.
Set~E is the normalized class (iii) rooms with $12\le W\le 20$ and
$n_R\le 2r$, with marks not both in $\Z^2$ that are admissible for one of the
$10$ lattices of index $2$ or $4$, plus the index-$6$ lattices for the one
$W=16$, area-$8$ room. Pairs in $\frac12\Z^2$ are skipped at area at most $14$,
because Part~A already lit them (d50 \S7). For each pair, the checker shoots
billiard paths in the four-fold unfolding with $p+q\le 40$, and then up to
$120$. An independent verifier \src{d56-final-check/src/verify_plane.c}
re-checks every lighting path and rejects any path through a corner. A lit pair
needs no certificate test: by Theorem~\ref{thm:U} it is not certified.
Survivors and a deterministic $1\%$ sample go to a certificate fallback.

\subsection{Counts}

\begin{table}[h]
\centering
\small
\begin{tabular}{@{}lrrr@{}}
\toprule
Run & Rooms & Pairs & CPU seconds \\
\midrule
d46 full, $W=20$, areas allowed by $F^*$ & 13,286,070,297 & 52,725,865,264 & 87,830 \\
d46 pilot, $W=14$, $16$, $18$ & 471,475,261 & 1,148,517,174 & 1,606 \\
d46 recheck, $W\le 20$, area $\le 20$, $F^*$ off & 427,403,072 & 135,236,864 & 1,094 \\
d56 Set A (D20) & 169,672,127 & 285,299,577 & 4,095 \\
d56 Set E (non-lattice marks) & 945,884 & 5,238,044,332 & 12,845 \\
\midrule
Total & 14,355,566,641 & 59,532,963,211 & \\
\bottomrule
\end{tabular}
\caption{All runs, with $0$ positives, $0$ survivors, and $0$ certified pairs.
Without the recheck row the pair total is $59{,}397{,}726{,}347$, the figure
behind ``at least $59.4$ billion'' in the results record. Sources:
\protect\src{d46-cert-check/runs/*/aggregate.md},
\protect\src{d56-final-check/runs/full/aggregate.md}.}
\label{tab:counts}
\end{table}

The pilot splits as $14{,}686$ rooms and $64{,}921$ pairs at $W=14$;
$12{,}231{,}060$ and $37{,}817{,}546$ at $W=16$; and $459{,}229{,}515$ and
$1{,}110{,}634{,}707$ at $W=18$. In Set~A all $285{,}299{,}577$ pairs are lit
at bound $40$. In Set~E, $5{,}236{,}565{,}128$ pairs are lit at bound $40$ and
$1{,}479{,}204$ need bound $120$. Set~E also ran an index-$3$ regression on
$9{,}631$ rooms ($43{,}623{,}736$ pairs) and $8{,}296$ index-$6$ pairs, all lit.
The certificate sample tested $55{,}231{,}932$ pairs ($2{,}852{,}044$ in Set~A
and $52{,}379{,}888$ in Set~E): $0$ were fully certified. Set~E recorded
$3{,}454{,}268$ single-target $\Par$ instances and $16$ $\Cut$ instances, none
completing a lift pair. $42{,}693{,}277$ sampled pairs were admissible only for
a sheared lattice, which the certificate code does not implement; they are lit,
so this does not affect the result.

\subsection{Independent checks}
\begin{itemize}
\item \emph{File checks.} For d46, all $500$ job files are present and all
$5{,}500$ per-cell lines read ``positives 0''. For d56, $920$ summary lines read
OK and the sums give $0$ survivors, $0$ certified pairs, $0$ tripwires, and
$0$ verification failures. The d56 \texttt{aggregate.py} crashed in its
reconciliation step; the coordinator fixed an operator-precedence bug and
reran the aggregation. No job was rerun.
\item \emph{Counts.} The d46 recheck reproduces the d9 census room and pair
counts at $W\le 20$, area at most $20$, and finds $0$ positives in
$427{,}403{,}072$ rooms. Set~A's room count equals the census count
$169{,}672{,}127$ in total and per job residue. Set~E's room counts match the
d50 table for each $W$: $22$, $287$, $3{,}875$, $59{,}893$, $881{,}807$.
\item \emph{Filter.} With and without $F^*$ at $W=14$, all
$54{,}305{,}681$ rooms ($126{,}988{,}441$ pairs) give $0$ positives
(\src{d46-cert-check/runs/val_W14_nofilter/aggregate.md}).
\item \emph{Shooting.} The d56 shooter agrees with the exact decider
\texttt{orbitc} on $32{,}127$ pairs, including the dark pair of $\Rr$
(\src{d56-final-check/README.md}).
\item \emph{Lemma tripwires.} d56 tests each sampled certificate against
Lemmas~\ref{lem:Z}, \ref{lem:Zprime}, \ref{lem:R0} and the D20 shape
$(4,5,2)$. All tripwire counts are $0$.
\item \emph{Reruns for this manuscript.} We reran
\texttt{d46-cert-check/aggregate.py} on the full run and the d45 and d62 check
scripts; the outputs match the stored records. Section~\ref{sec:repro} lists
the commands.
\end{itemize}

\subsection{Controls}\label{sec:controls}
The 22-wall room $\Rr$ has $W=22$, $r=9$, area $15$, and one dark pair,
$(0,0)$ and $(2,2)$ in the coordinates of \cite{Q22}. In the d9 census
coordinates its cells are $(0,2)$, $(0,3)$, $(1,1)$, $(1,2)$, $(1,3)$,
$(1,4)$, $(2,0)$, $(2,1)$, $(2,4)$, $(2,5)$, $(3,4)$, $(3,5)$, $(4,3)$,
$(4,4)$, $(5,3)$, and the pair is $S=(1,3)$, $T=(3,5)$.
\begin{itemize}
\item \emph{Certified.} $B_b$, $B_c$, $B_d$ have $5$, $15$, $5$ blocks and
separate all four lift pairs; $(\delta_b,\delta_c,\delta_d)=(3,1,3)$. So the
pair is certified dark by $\Par$ at all three targets over $B$.
\item \emph{d46 pipeline.} An end-to-end run at $r=9$, area $15$, with $F^*$ on
checks $556{,}286$ rooms and $3{,}789$ pairs and finds exactly one positive:
$S=(1,3)$, $T=(3,5)$, blocks $5,15,5$
(\src{d46-cert-check/runs/e2e_r9_a15/aggregate.md}
and \src{d46-cert-check/positive_e2e.txt}). The census-range run at $W=22$, area
at most $16$, finds the same single positive in $3{,}326{,}713$ rooms. $F^*$
passes every room at $r=9$, $n=15$, so this run tests the pipeline but not the
pruning.
\item \emph{d56 checker.} On $\Rr$ the pair survives shooting to bound $120$,
$\Par$ certifies it at $b$, $c$, $d$ for all four lift pairs ($12$
instances), and the checker exits with its positive code.
\item \emph{Second proof.} At $r=9$ the fold enumeration of
Section~\ref{sec:second} finds three positive covers, all equal to $\Rr$ up to
lattice isometry.
\item \emph{Negative control.} The 12-wall cross
$([0,4]\times[1,3])\cup([1,3]\times[0,4])$ with $S=(1,2)$, $T=(3,2)$ has $\Par$
at $d$, $\Cut$ at $c$, and nothing at $b$: it is lit and uncertified.
\item \emph{Second dark room.} The room P6 ($W=28$, area $20$) is positive at
$S=(1,0)$, $T=(3,-2)$ with blocks $5,5,5$. It lies outside the enumerator's
range and was checked by direct input
(\src{d46-cert-check/val_named.py}).
\end{itemize}

\section{A second proof by fold enumeration}\label{sec:second}

The d52 memo proves the core case again without enumerating rooms. The d54
referee re-checked it, completed its validation, and removed one input.

\begin{theorem}[d52, as signed in d54; computer-assisted]\label{thm:second}
(1) Let $P$ be a strictly simple orthogonal lattice room with $r\le 8$ whose
corners lie in at least $3$ parity classes. Assume Theorem~\ref{thm:U},
Proposition~\ref{prop:R3}, Lemma~\ref{lem:factL}, d36.3, and the d39/d43
reduction (Proposition~\ref{prop:F1}, which uses the d9 census for essential
$\Rig$). Then no two distinct interior lattice points of one class are dark by
$\Par$, $\Cut$, or $\Rig$ certificates over the unit pillowcase. This holds at
every area.

(2) Under the same inputs except the d39/d43 $\Rig$ reduction: any strictly
simple room with $r=9$, corners in at least $3$ classes, and a same-class
interior pair certified dark by three $\Par$ certificates is $\Rr$ up to
lattice isometry, with its one pair.
\end{theorem}

\paragraph{Method.}
Write $K_x=N_x/\iota$. Case~A: some target class $x\ne a$ has
$\delta_x\ge 3$. Fix $x=(0,0)$ by translation and enumerate the small
quotient $K_x$ by a low-index search for transitive sets under the four edge
reflections. The constraints are Lemma~\ref{lem:factL}, the seam conditions,
at most one interior $\pi$-cone, $\delta z\le r$ with $\delta$ even when there
is a cone, singular points in at least $3$ classes, and two non-reflex
class-$a$ points. Then enumerate every degree-$\delta$ cover of $K_x$ by one
involution per boundary arc between convex corners and a fixed-point-free
involution across a slit from the cone. For each cover, recompute $B_y$ for all
four classes on the cover itself and test separation.
Case~B: all three targets have $\delta\le 2$. Each quotient is then Galois,
d36.3 makes them a chain, and the coarsest $B_*$ separates, has $\delta_*\le 2$,
and has no regular point over any target. Summing over the three targets gives
$3d\le 4z+4$, so $n\le 12$.

\paragraph{Results.}
At $r\le 8$, Case~A has $106$ admissible quotient shapes ($56$ up to mirror),
with $d$ from $3$ to $9$ and $z\in\{1,2\}$; $63$ have an interior cone. Over
$185$ (shape, $\delta$) cases the search finds $4{,}710$ covers per case,
$4{,}541$ with corners in at least $3$ classes; these are $4{,}420$ and
$4{,}275$ distinct W-sets, $1{,}744$ discs up to lattice isometry, and $1{,}113$
strictly simple rooms. No mark pair is separated in all three classes; the best
is $2$ of $3$ ($8$ pairs). Case~B (d54) has $190$ quotient shapes, $43$
in-scope degree-$1$ rooms and $23$ in-scope degree-$2$ covers, and $0$
positives. At $r=9$, Case~A covers $12{,}488$ discs and finds $3$ positive
covers, all equal to $\Rr$; Case~B at $r=9$ finds $0$ positives.

\paragraph{Validation.}
A second, model-free cover enumerator agrees with the arc model on $185$ of
$185$ cases ($10$ of them finished by d54). A brute-force pass over all
$13{,}876$ immersed W-sets with $n\le 9$ recovers all $24$ incidences; d54
notes this test exercises only $d=3$. Every $\delta\ge 3$, $d\ge 2$ quotient of
a found cover ($4{,}018$) lies in the shape list, and every cover with
$5\le r\le 8$, $n\ge 21$ passes $F^*$ ($2{,}845$ of $2{,}845$). Two
independent separation codes agree on all $4{,}275$ in-scope covers.

\paragraph{Independence.}
The code shares nothing with d44 or d46, and the method enumerates quotients,
not rooms. With Case~B, the proof needs neither $F^*$ nor d2 claim~4, which
rests on the computer-assisted theorem of \cite{K6}. It still shares with d46:
Theorem~\ref{thm:U}, Proposition~\ref{prop:R3}, Lemma~\ref{lem:factL}, the
d39/d43 reduction with the d9 census for essential $\Rig$, the bound
$\delta z\le r$, and the scope (same-class lattice marks, unit pillowcase, at
least $3$ corner classes). The no-positive part also holds for immersed,
slit, and pinched discs with room-local types.

\paragraph{Not covered.}
Theorem~\ref{thm:second} does not treat other pillowcase bases, different-class
marks, or non-lattice marks; for those, Theorem~\ref{thm:main} uses d47 to d50
and d56. At $r=9$ it does not treat rooms certified with essential $\Rig$, so it
does not prove that $\Rr$ is the only certified-dark room with $22$ walls. The
d54 referee did not rerun d52's own enumerations; the counts above are quoted
from d52's logs and d54's checks.

\section{The fold converse}\label{sec:converse}

Theorem~\ref{thm:main} would decide the $W\le 20$ case of the illumination
question if every dark pair were certified. That is conjecture $U^*$ of d10
\S2.4. The d61 memo studies the remaining step one target at a time.

\paragraph{States.}
A \emph{state} is a basis in the pure-braid orbit of the monodromy tuple. Each
state gives a tail arc $\alpha\colon a\to x$. Every isotopy class of arc
$a\to x$ is a tail arc in some state, because the pure braid group acts on
slopes through $\Gamma(2)$ (d62 \S1). A pair is \emph{$x$-lit} when some lift
pair is joined along some tail arc to $x$, and \emph{$x$-dark} otherwise.

\begin{lemma}[Lemma FC; d61, reviewed in d62]\label{lem:FC}
Fix a state and an arc $\alpha\colon a\to x$. Let $\tilde S\ne\tilde T$ be regular
points of $Q$ over $a$ with the same image $u$ in a factor $N$, and let
$K=K_N(\alpha,u)$.
\begin{enumerate}
\item $K$ is a path or a cycle. If it is a cycle, every vertex is
$N$-regular.
\item \emph{Cycle case.} Let $\mu$ be the monodromy of $Q\to N$ around $K$,
acting on the fibre over $u$. Then $\tilde S$ and $\tilde T$ are joined along
$\alpha$ exactly when $\tilde T$ lies in the $\langle\mu\rangle$-orbit of $\tilde S$.
\item \emph{Path case.} For each live-pole end $p$ of $K$, let $\tau_p$ be the
local monodromy of $Q\to N$ at $p$, transported to the block of $u$ along $K$
and restricted to its $2$-cycles. If $u$ is a pole, $\tau_u$ is $g_a$ on the
block of $u$, restricted to $2$-cycles. These maps act on the sheets in the
block of $u$. Then $\tilde S$ and $\tilde T$ are joined along $\alpha$ exactly when
an alternating word in the at most two $\tau_p$ carries some sheet of $\tilde S$
to some sheet of $\tilde T$.
\item The image in $N$ of a light path is a closed walk at $u$ inside $K$ that
reverses only at live-pole ends and never meets a point over an $N$-zero.
\item If $\mu$ moves $\tilde S$, both discs of $N$ bounded by $K$ contain a pole
over which $Q\to N$ is branched.
\item Replacing $\tau_p$ by the full local monodromy changes the verdict only
through its $3$-cycles, which are reflex corners.
\end{enumerate}
\end{lemma}

The statement on sheets in part~3 is the d62 referee's fix G1: when $u$ is a
pole, the two sheets of a regular point lie in one block, and the maps do not
descend to points. The proof uses Lemma~\ref{lem:factL} and genus $0$ only; $N$
need not be minimal. Over an $N$-regular vertex, $Q$ has $e=1$ and the image
walk passes straight through; over a live pole only an $e=2$ point continues,
and the image walk reverses. Part~5 is the last step of the $\Cut$ proof in
Theorem~\ref{thm:U}.

\paragraph{Consequences.}
So a pair is $x$-lit exactly when, in some state, the partial dihedral or
cyclic action of Lemma~\ref{lem:FC} joins two of its lifts. No group replaces
this family: the monodromy group of $Q\to N$ is transitive on the fibre, and
using full local permutations, which ignores corner avoidance, lights the dark
pairs of $\Rr$ and P6 (d61 \S3).

\begin{proposition}[Degree bound; d61 \S6, reviewed in d62]\label{prop:degN}
$N_x$ has no regular point over $x$, and
$\deg N_x\le (r+4)+3r_x\le 4r+4\le 36$ at $r\le 8$, where $r_x$ counts reflex
corners of class $x$.
\end{proposition}
\begin{proof}
Points of $N_x$ over $x$ have length $1$ or $3$, so $\deg N_x=P_x+3Z_x$.
Gauss--Bonnet gives $P-Z=4$, and Lemma~\ref{lem:factL} gives $Z\le r$ and
$Z_x\le r_x$.
\end{proof}

\paragraph{The open direction.}
By Proposition~\ref{prop:R3}, ``$x$-dark implies a certificate at $x$'' is
equivalent to the following \emph{realization lemma} (d61 \S5, with the d62
fix G3 to its quantifiers): let $\tilde S\ne\tilde T$ be regular $a$-points with
one image $u$ in $N_x$; if $\Rig(x)$ fails over $DP/R_y$ for both
$y\notin\{a,x\}$ and $\Cut(x)$ holds for no factor, then in some state the
$2$-cycles of $g_a$ and $g_x$ inside the blocks of $u$ connect $\tilde S$ to
$\tilde T$. The realization lemma is \textbf{open}. Proposition~\ref{prop:degN}
makes $N_x$ range over a finite set, but $Q\to N_x$ has unbounded degree, so
the lemma is not a finite check. If it were proved for orthogonal lattice rooms,
Theorem~\ref{thm:main} and the rational-to-real reduction \cite{RR} would give
illumination of every pair at $W\le 20$ (d61 \S6).

\paragraph{Tests.}
On the four named rooms (cross12, dumb12, $\Rr$, P6), d61 compares the lemma's
verdict with a direct walk on $405$ rows and $992{,}988$ (state, arc, lift pair)
instances, with $0$ disagreements. The d62 referee found that d61's verdict is
a segmented re-simulation of the same walk, so these counts test
Lemma~\ref{lem:factL} and Lemma~\ref{lem:proj} more than the algebraic form
(finding G2). The referee's own implementation of the stated form agrees with a
direct walk on $969{,}372$ instances, and an exact integer billiard program
agrees with d61's rows at $N=B$ with $0$ mismatches. Every $x$-dark row in the
tests carries a certificate: $147$ by $\Par$ and $56$ by $\Cut$.

\section{Fold mechanism needs \texorpdfstring{$r\ge 9$}{r >= 9}}\label{sec:r9}

This section records a later result, written on 20 November 2026 in
\src{research/qclass-closure/bbbb/fold-converse.md} and reviewed in d105
(verdict: sound, with one fixable gap). It is separate from
Theorem~\ref{thm:main} and uses its own inputs.

\paragraph{Setting.}
Let $P$ be a strictly simple orthogonal lattice room, $Q=DP$, and $S\ne T$
interior lattice points of one class $a$, with lifts $s,\iota s$ and
$t,\iota t$. An \emph{R22-type parity fold} is a factor
$f\colon Q\to N\to B$ of inner degree $k\ge 2$ and outer degree $m>1$ such that:
all four lift images are regular points over $a$; the image sets
$\{f(s),f(\iota s)\}$ and $\{f(t),f(\iota t)\}$ are disjoint; and $N$ has no
regular point over any of the three off-classes $b,c,d$. Such a factor is a
$\Par$ certificate at all three targets, so the pair is dark by
Theorem~\ref{thm:U}.

\begin{proposition}[Fold mechanism needs $r\ge 9$; d105, sound with fix]\label{prop:r9}
If $P$ has an R22-type parity fold for some interior pair, then $r\ge 9$.
The bound is attained by $\Rr$, with $(k,m,z)=(3,5,3)$.
\end{proposition}

The separation hypothesis is the d105 fix. The source's main hypothesis
asked only for two distinct regular marked images; d105 shows the inner-degree-two
step needs separation, and notes that the mechanism supplies it.

\begin{proof}[Proof (from the source, with the d105 fix)]
\emph{Counting.} $N\to B$ has local degrees in $\{1,2,3\}$ and $N$ is a sphere.
With $q$ regular points and $z$ zeros, Riemann--Hurwitz gives $q+2z=2m-2$, and
the poles number $z+4$. All $q$ regular points lie over $a$ and use two sheets
each, and the marks need $q\ge 2$. So $2\le q\le m/2$ and
$z+2\le m\le 4(z+1)/3$. Over a zero of $N$ every preimage has $e=1$, so each
zero has $k$ reflex preimages; with $c_j$ points of local degree $j$ over
poles, $r=kz+c_3$ and $c_2+2c_3=2k-2$. At $r\le 8$ the only possibilities are
$(z,m,k)\in\{(2,4,2),(2,4,3),(2,4,4),(3,5,2),(4,6,2)\}$, each with $q=2$. The
d105 script \src{d105-review-fold/step2.py} confirms this list.

\emph{Refinement.} Replace the partition by
$\Pi'=B_b\vee B_c\vee B_d$, the finest monodromy-invariant equivalence that
merges every $2$-cycle at the three off-classes. It refines the given
partition, its quotient still has no regular off-class point, the marked images
stay regular and separated, and its outer degree is at least $m$. It is
invariant under the edge involutions by Proposition~\ref{prop:R3}(3).

\emph{Discrete case.} If $\Pi'$ is discrete, the three off-classes are
corner-only in $P$. Each of the four cardinal neighbours of $S$ and of $T$ has
two cells of the mark's star, so it is a reflex corner. A coincidence of a
neighbour of $S$ with one of $T$ would put an interior point at a corner-only
class, so the eight are distinct (d105 corrects the source here: the case is
vacuous, not ``directly visible''). A diagonal neighbour of $S$ in class $d$
is a ninth reflex corner (d14 Lemma R6). So $r\ge 9$.

\emph{Inner degree two.} In the rows with $k=2$, $q=2$ and separation put both
lifts of each mark over one quotient point. The deck involution exchanges the
halves at both marks; by the axis lemma of d91 and d99 (\src{research/qclass-closure/bbbb/bbbb-split.md}),
it is an axial reflection of the room fixing both marks. The interior of a
simply connected room symmetric in a line meets the line in one interval, so
the segment $ST$ is interior and the pair is lit. This contradicts darkness.

\emph{Outer degree four.} In the row $m=4$, $z=2$, label the four blocks. At
$a$ the corner permutation has type $(2,2)$; the two zeros give type $(3,1)$ at
two off-classes, and the third off-class has the identity. All four corner
permutations are even, so the four edge involutions have one sign; even
involutions lie in $V_4$, which has no $3$-cycle, so all four are
transpositions. At the mark over the $a$-cycle $(12)$, the two eastern cells
share one label in $\{1,2\}$ and the two western cells the other. A $3$-cycle
at the horizontally adjacent class contains exactly one of $1,2$, so on that
side two cells with one label meet at a reflex corner, whose upstairs cycle
must carry three distinct labels. This contradiction forces the identity at the
horizontally adjacent class; the other mark forces it at the vertically
adjacent class. Only one off-class has the identity, a contradiction. The d105
script finds $72$ tuples in $S_4$ with the required cycle types, none with the
identity at both adjacent classes.
\end{proof}

\begin{remark}[Census route]
The source also gives a shorter route: in every row of the table the upstairs
area $km$ is at most $16$, and the d9 census excludes dark same-class interior
pairs at $r\le 8$ and area at most $20$. The source also eliminates the
$z=2$, $m=4$ row for a physical polygon base by a direct check of the two
tetrominoes with two reflex corners; d105 reproduces it.
\end{remark}

\begin{proposition}[A proper fold is not necessary; source Step 1, confirmed in d105]\label{prop:R34}
The strictly simple room R34-base, with $19$ cells, $19$ convex and $15$ reflex
corners ($W=34$), interior points $(2,0)$, $(6,0)$, and smooth lattice points
$(4,0)$, $(8,-2)$, has a dark pair, and $\deg(Q\to B)=19$ is prime. So its
darkness comes from no factor with $k\ge 2$ and $m\ge 2$.
\end{proposition}
\begin{proof}[Proof summary]
All four regular lattice sites are even--even. For any proposed shot between
the two interior points, the unfolded displacement is even--even and nonzero,
so its first primitive lattice step lies strictly before the target and is not
even--even. Reflection in integer lines preserves parity, so the folded image
is a corner. Since $19$ is prime, the monodromy group has no proper block
system. This is a parity argument in the sense of Tokarsky \cite{Tokarsky}.
The d105 script \src{d105-review-fold/step1.py} reproduces the star count.
\end{proof}

\begin{remark}[Scope of Proposition~\ref{prop:r9}]
In a room where every class quotient $N_x$ is $B$ (type BBBB), a factor with
no regular point over an off-class is coarser than $B_x$, so $m=1$. So
Proposition~\ref{prop:r9} is vacuous on BBBB rooms at every $r$, and it rules
out nothing in the open core. At $r\le 8$ it reproves, without the census, a
part of a case that the Q-class closure (d102) and d64 already close (d105).
\end{remark}

\section{Scope and limitations}\label{sec:limits}

\paragraph{Computer assistance.}
Theorem~\ref{thm:main} is computer-assisted. Its finite part runs C and Python
programs that are not formally verified. The full $W=20$ room check (d46) has
one implementation; the coordinator checked its output files but did not rerun
it. Theorem~\ref{thm:second} reproduces the core case with independent code,
and d54 did not rerun d52's own enumerations. The d56 shooter is validated
against an exact decider on $32{,}127$ pairs, not on every pair. No d56 test
produced a positive $\Rig$ separation, so its $\Rig$ code is checked against a
reference only in the ``not separated'' direction.

\paragraph{Review.}
Every theory input has one independent AI referee, usually on the same model
family as the author. No human mathematician has refereed any part. The
referees' repairs in d43, d48, d50, d54, d62, and d105 are themselves
unreviewed. Some inputs carry recorded author confidence below certainty: d47
gave $75\%$ for the transfer lemma and for Proposition~\ref{prop:diffclass}
before review. The d9 census, which Propositions~\ref{prop:F1} and
\ref{prop:even} use, has computer re-checks but no line-by-line review. The
parity-factor input d2 claim~4 rests on the computer-assisted theorem of
\cite{K6}; Theorem~\ref{thm:second} avoids it, but $F^*$ in
Theorem~\ref{thm:main} does not.

\paragraph{Gaps flagged in the sources.}
\begin{itemize}
\item The d62 referee did not re-verify that d46, d52, and d56 allow the
certifying factor to change with the target and the lift pair (finding C4).
Proposition~\ref{prop:F1} allows this change; we found no separate check of it
in the d46 or d56 code records.
\item The completeness argument for the d56 D20 shortcut
(\src{d56-final-check/shortcut/NOTES.md}) is not refereed. It is not load-bearing:
the full Set~A run lights every D20 pair.
\item The d48 treatment of $(2,4,4)$ bases and the witness-route bounds are
sketches.
\item Lemma~\ref{lem:FC}: the d62 fixes G1 and G3 are not written into d61.
\item Proposition~\ref{prop:r9} uses the d91/d99 axis lemma and d14 Lemma R6,
which we cite without restating their proofs.
\end{itemize}

\paragraph{What the theorem does not cover.}
\begin{itemize}
\item \emph{The lit side.} Theorem~\ref{thm:main} shows that no pair is
certified dark. It does not show that uncertified pairs are lit. That is
conjecture $U^*$, or per class the realization lemma of
Section~\ref{sec:converse}, and it is open.
\item \emph{Other mechanisms.} Only $\Par$, $\Cut$, and $\Rig$ over pillowcase
bases are covered. A dark pair certified by some other mechanism is not
excluded.
\item \emph{Irrational pairs.} Certificates need rational marks. The
rational-to-real reduction \cite{RR} transfers illumination, not certificates.
\item \emph{Boundary marks.} The marks are interior points. Statements such as
R4 are false for smooth wall marks.
\item \emph{45-degree walls.} Rooms with $45$-degree walls are not covered; d51
(unreviewed) finds no comparable bound.
\item \emph{$W\ge 22$.} $\Rr$ shows the statement fails at $W=22$.
\end{itemize}

\paragraph{Discrepancies between sources.}
\begin{itemize}
\item The results record lists the pair count of the d46 recheck as ``not
stated''; the run aggregate gives $135{,}236{,}864$. Including it raises the
total from $59{,}397{,}726{,}347$ to $59{,}532{,}963{,}211$.
\item d52 reports $4{,}710$ covers; these are counted per (shape, $\delta$)
case with repeats. d54 counts $4{,}420$ distinct W-sets ($4{,}275$ in scope).
\item d52 states the $r=9$ classification under ``some class has
$\delta\ge 3$''; d54 requires a target class $x\ne a$.
\item d47 Proposition~4.1 lists $d\in\{3,4,5\}$; d48 shows $d=5$.
\item d49 says ``$\Cut$ forces $d=3$'' and counts $2{,}052{,}226$ rooms for
its finite check; d50 corrects the first to ``with a zero'' and the room set to
$945{,}884$.
\item The source of Section~\ref{sec:r9} omits the separation hypothesis and
calls the discrete coincidence case ``directly visible''; d105 corrects both.
\item An older status file tags d61's lit direction as unreviewed; d62
records the review.
\end{itemize}

\section{Reproduction}\label{sec:repro}

The repository mirrors the research directory at
\[\text{\src{research/illumination-generalization-2026-11-06/fable/directions-2026-11-09/}}\]
Below, paths that begin with \texttt{research/} are relative to the
repository root, and all other source paths are relative to that directory.
Some original logs and the d46 README are not in the mirror; the run
aggregates and scripts cited here are. To run a script that expects the
original layout, use the restoration command in \src{docs/REPRODUCIBILITY.md}.

\begin{center}
\small
\begin{tabular}{@{}>{\raggedright\arraybackslash}p{0.2\textwidth}>{\raggedright\arraybackslash}p{0.6\textwidth}>{\raggedright\arraybackslash}p{0.13\textwidth}@{}}
\toprule
Result & Command (from the named directory) & Cost \\
\midrule
Room check (F3) & In \texttt{d46-cert-check/}: \texttt{./run\_full.sh 8 24}; then \texttt{python3 aggregate.py runs/full\_W20 --cells} & 24.4 CPU-h \\
Residual checks (F5) & In \texttt{d56-final-check/}: \texttt{./run\_full.sh 8 10}; the script compiles \texttt{check} from \texttt{src/} when no binary is present & 4.7 CPU-h \\
Second proof & In \texttt{d52-fold-enumeration/}: \texttt{python3 -B enum\_K.py 2 12 data/K\_z2\_d12.json}, \texttt{analyze\_K.py data/K\_z2\_d12.json 8 2 data/K\_adm\_r8.json}, \texttt{enum\_P.py data/K\_adm\_r8.json 8 data/P\_r8.json} & about 30 CPU-min \\
Case B & In \texttt{d54-review-d52/}: \texttt{python3 -B case\_b.py 8} & 21 CPU-s \\
$F^*$ check & In \texttt{d45-review-fstar/}: \texttt{python3 -B check\_fstar.py} & under 1 CPU-s \\
Lemma FC & \texttt{d61-fold-converse/fold\_converse.py}; \texttt{d62-review-d61/check\_fc.py ROOM CAP}, \texttt{billiard.py} & about 3 CPU-min \\
Degree bound & In \texttt{d62-review-d61/}: \texttt{python3 -B deg\_bound.py} & under 1 CPU-s \\
$r\ge 9$, R34 & In \texttt{d105-review-fold/}: \texttt{python3 -B step1.py}, \texttt{python3 -B step2.py} & about 2 CPU-s \\
\bottomrule
\end{tabular}
\end{center}

\begin{sloppypar}
For this manuscript we reran, in a temporary copy of the original layout,
\texttt{aggregate.py} on
\texttt{runs/full\_W20}, \texttt{check\_fstar.py}, \texttt{deg\_bound.py},
\texttt{step1.py}, \texttt{step2.py}, and \texttt{case\_b.py 8}. Each output
matches its stored record, apart from CPU times and one log header. We did not
rerun the d46, d52, or d56 enumerations.
\end{sloppypar}

The main sources are listed below; \texttt{provenance.json} records their
SHA-256 values. The results record is cited for \S\S2.3--2.7, 3.3--3.7, 5,
10, and 11.
\begin{itemize}\setlength{\itemsep}{0pt}\small
\item \src{d70-writeup/RESULTS.md}
\item \src{d10-mechanisms/theory.md}
\item \src{d2-review/review.md}
\item \src{d36-amalgamation/astra.md}
\item \src{d39-mixed-class/astra.md}
\item \src{d43-review-d39/review.md}
\item \src{d44-cert-check-sizing/README.md}
\item \src{d45-review-fstar/review.md}
\item \src{d46-cert-check/runs/full_W20/aggregate.md}
\item \src{d47-cert-gaps/opus.md}
\item \src{d48-review-d47/review.md}
\item \src{d49-even-index/memo.md}
\item \src{d50-review-d49/review.md}
\item \src{d52-fold-enumeration/memo.md}
\item \src{d54-review-d52/review.md}
\item \src{d56-final-check/README.md}
\item \src{d56-final-check/runs/full/aggregate.md}
\item \src{d61-fold-converse/README.md}
\item \src{d62-review-d61/review.md}
\item \src{d105-review-fold/review.md}
\item \src{d9-census-w22/README.md}
\item \src{d9-census-w22/partA-README.md}
\item \src{coordinator-log-archive-1.md}
\item \src{research/qclass-closure/bbbb/fold-converse.md}
\end{itemize}

\paragraph{Acknowledgments and AI assistance.}
The authors used ChatGPT and Claude to develop and check arguments, write and
run the programs, and prepare the manuscript. The authors take responsibility
for its content.

\begin{thebibliography}{9}
\bibitem{Tokarsky}
G.~W. Tokarsky, ``Polygonal Rooms Not Illuminable from Every Point,''
\emph{The American Mathematical Monthly} \textbf{102} (1995), no.~10, 867--879.
\url{https://doi.org/10.1080/00029890.1995.12004679}.

\bibitem{Q22}
P.~Wang and D.~Wang, \emph{A Simple Orthogonal 22-Gon with Exactly One
Non-Illuminating Pair}, working draft v0.7 (2026), \texttt{preprints/twenty-two-wall/} in this repository.

\bibitem{RR}
P.~Wang and D.~Wang, \emph{Rational Pairs Suffice for Illumination in
Orthogonal Lattice Rooms}, working draft v0.1 (2026),
\texttt{preprints/real-to-rational/} in this repository.

\bibitem{K6}
D.~Wang and P.~Wang, \emph{A restricted 22-wall minimum and a parity-fold
corollary}, working draft v0.1 (2026), \texttt{preprints/restricted-minimum-k6/} in this repository,
Theorem~1.1.
\end{thebibliography}
\end{document}
